% Part: many-valued-logic
% Chapter: infinite-valued-logics
% Section: lukasiewicz

\documentclass[../../../include/open-logic-section]{subfiles}

\begin{document}

\olfileid{mvl}{inf}{luk}

\olsection{منطق~\L ukasiewicz}

\begin{editorial}
  ليس هذا إلا «بذرة» وجيزة لقسم في منطق~\L ukasiewicz اللانهائي القيم.
\end{editorial}

\begin{defn}\ollabel{def:lukasiewicz} مصفوفة منطق~\L ukasiewicz
اللانهائي القيم~$\LogLuk[\infty]$ هي:
\begin{enumerate}
  \item الجزء الخالي من الثوابت~$\Lang L_0^{-}$ من لغة القضايا القياسية، ذو الروابط
  $\lnot$ و$\land$ و$\lor$ و$\lif$.
  % تصحيح مصدر مشترك (MVL-L0-I1): حُددت اللغة الخالية من lfalse التي تفسر المصفوفة جميع روابطها.
  \item مجموعة قيم الصدق~$V_\infty$.
  \item القيمة المميزة وحدها هي~$1$؛ أي~$V^+ = \{1\}$.
  \item دوال الصدق هي الدوال الآتية:
  \begin{align*}
    \tf{\lnot}[\LogLuk](x) & = 1 - x\\
    \tf{\land}[\LogLuk](x,y) & = \min(x,y)\\
    \tf{\lor}[\LogLuk](x,y) & = \max(x,y)\\
    \tf{\lif}[\LogLuk](x,y) & = \min(1,1-(x-y)) = \begin{cases}
      1 & \text{إذا كان } x \le y\\
      1-(x-y) & \text{في غير ذلك.}
    \end{cases}
    \end{align*}
\end{enumerate}
وأما منطق~\L ukasiewicz ذو~$m$ من القيم فيعرف على الوجه نفسه، إلا
أن~$V = V_m$.
\end{defn}

\begin{prop}
  المنطق~$\LogLuk[3]$ المعرف في \olref[thr][luk]{def:lukasiewicz}
  هو بعينه~$\LogLuk[3]$ المعرف في \olref{def:lukasiewicz}.
\end{prop}

\begin{proof}
  يتبين ذلك إذا قابلنا جداول الروابط في
  \olref[thr][luk]{def:lukasiewicz} بجداول الصدق المستخرجة من
  معادلات \olref{def:lukasiewicz}:
  \begin{center}
    \begin{tabular}{c|c}
      $\tf{\lnot}$ & \\
      \hline
      $1$ & $0$ \\
      $1/2$ & $1/2$ \\
      $0$ & $1$
    \end{tabular}
    \quad
    \begin{tabular}{c|ccc}
      $\tf{\land}[\LogLuk[3]]$ & $1$ & $1/2$ & $0$ \\
      \hline
      $1$ & $1$ & $1/2$ & $0$ \\
      $1/2$ & $1/2$ & $1/2$ & $0$\\
      $0$ & $0$ & $0$ & $0$
    \end{tabular}
    \\[2ex]
    \begin{tabular}{c|ccc}
      $\tf{\lor}[\LogLuk[3]]$ & $1$ & $1/2$ & $0$ \\
      \hline
      $1$ & $1$ & $1$ & $1$ \\
      $1/2$ & $1$ & $1/2$ & $1/2$ \\
      $0$ & $1$ & $1/2$ & $0$
    \end{tabular}
    \quad
    \begin{tabular}{c|ccc}
      $\tf{\lif}[\LogLuk[3]]$ & $1$ & $1/2$ & $0$ \\
      \hline
      $1$ & $1$ & $1/2$ & $0$ \\
      $1/2$ & $1$ & $1$ & $1/2$  \\
      $0$ & $1$ & $1$ & $1$
    \end{tabular}
  \end{center}
\end{proof}

\begin{prop}\ollabel{prop:luk-infty-m}
  إذا كان~$\Gamma \Entails[\LogLuk[\infty]] !B$، كان
  $\Gamma \Entails[\LogLuk[m]] !B$ لكل~$m \ge 2$.
\end{prop}

\begin{proof}
  إثبات ذلك تمرين.
\end{proof}

\begin{prob}
  أثبت \olref[mvl][inf][luk]{prop:luk-infty-m}.
\end{prob}

% تصحيح مصدر (0400-I1): لا يصح إطلاق العكس لمجموعة مقدمات غير منتهية.
وأما عكس هذا الحكم فيصح إذا كانت مجموعة المقدمات متناهية.

ومنطق~\L ukasiewicz اللانهائي القيم أشهر المنطقات الضبابية. وأما
الشرط في أدبيات المنطق الضبابي فكثيرًا ما يعرف بـ~$\lnot !A \lor !B$؛
والحاصل حينئذ منطق كليني قوي لانهائي القيم.

\begin{prob}
  بيّن أن~$(p \lif q) \lor (q \lif p)$ تحصيل منطقي في
  $\LogLuk[\infty]$.
\end{prob}

\end{document}
